%=====================================================================
%  Immigrant Owners and Trade Diversification
%  Draft paper: introduction, data, model, empirical strategy;
%  results, mechanisms, discussion and conclusion are outlined.
%  Compile: pdflatex (twice). Bibliography is embedded.
%=====================================================================
\documentclass[12pt]{article}

\usepackage[T1]{fontenc}
\usepackage[utf8]{inputenc}
\usepackage{mathpazo}
\usepackage[margin=1in]{geometry}
\usepackage{setspace}
\usepackage{amsmath,amssymb,amsthm}
\usepackage{booktabs,tabularx,array,multirow}
\usepackage{graphicx}
\usepackage{caption}
\usepackage{enumitem}
\usepackage[dvipsnames]{xcolor}
\usepackage[authoryear,round]{natbib}
\usepackage[colorlinks=true,linkcolor=MidnightBlue,citecolor=MidnightBlue,urlcolor=MidnightBlue]{hyperref}

\onehalfspacing
\setlength{\parskip}{0.35em}
\captionsetup{font=small,labelfont=bf}
\newcolumntype{Y}{>{\raggedright\arraybackslash}X}

% ---------- Theorem environments ----------
\newtheorem{proposition}{Proposition}
\newtheorem{lemma}{Lemma}
\newtheorem{assumption}{Assumption}
\theoremstyle{remark}
\newtheorem{remark}{Remark}

% ---------- Drafting macros ----------
% \result{...}  : placeholder for a finding to be filled in
% \verify{...}  : institutional/data detail to confirm before circulating
% \planned{...} : outline text for sections not yet written
\newcommand{\result}[1]{\textcolor{RoyalBlue}{[\textbf{Result:} #1]}}
\newcommand{\verify}[1]{\textcolor{BrickRed}{[\textbf{Verify:} #1]}}
\newcommand{\planned}[1]{\textcolor{OliveGreen}{#1}}
\newcommand{\figshell}[2]{%
  \begin{figure}[htbp]\centering
  \fbox{\begin{minipage}[c][1.9in][c]{0.85\textwidth}\centering
  \planned{\textit{Figure to be produced.}}\\[0.4em]
  \small\planned{#2}\end{minipage}}
  \caption{#1}\end{figure}}

% ---------- Notation ----------
\newcommand{\Post}{\mathrm{Post}}
\newcommand{\one}{\mathbf{1}}
\newcommand{\E}{\mathbb{E}}

\title{\textbf{Immigrant Owners and Trade Diversification}\thanks{%
Preliminary and incomplete; results to be added. Statistics Canada
disclosure statement and acknowledgements to be added.}}
\author{Nathanael [Surname]\\[0.2em]\normalsize [Department, University]}
\date{November 2026}

\begin{document}
\maketitle

\begin{abstract}
\noindent
Canada has long sought to diversify its international trade, but the
decisions that determine which countries and products a country trades are
made by firms. This paper asks whether the identity of a firm's owner is one
of those determinants. Using Canadian customs records linked to owner records,
I follow roughly 1,000 incumbent firms whose ownership passes from
non-immigrant to immigrant owners between 2005 and 2019, and compare each
firm's trade with its new owner's country of origin to its trade with other
countries, before and after the change. Because the business is held fixed
while the person in control changes, the design separates knowledge that moves
with the owner from the sorting of immigrants into firms already positioned to
trade with their origin countries. A model of trade-relationship choice
distinguishes three channels---lower fixed costs of forming relationships,
lower variable costs of running them, and a shift in demand toward
home-country goods---by their implications for the extensive and intensive
margins, for imports and exports, and for consumer and intermediate goods.
\result{Headline effect on trade with the origin country; the margin through
which it operates; whether it diversifies the firm's trade or reallocates it
from other partners.}

\medskip\noindent
\textbf{Keywords:} immigration; ownership; firm-level trade; trade
diversification; information frictions\\
\textbf{JEL:} F14, F22, L25, D83
\end{abstract}

\thispagestyle{empty}
\newpage
\setcounter{page}{1}

%=====================================================================
\section{Introduction}
\label{sec:intro}
%=====================================================================

Few objectives have proved as durable in Canadian trade policy as
diversification. Roughly three-quarters of Canada's merchandise exports go to
a single partner, the United States, and successive federal strategies have
aimed to expand trade with the rest of the world---an objective that has
acquired new urgency since the escalation of U.S. tariffs in 2025.%
\footnote{\verify{Exact U.S. share of Canadian merchandise exports for the
sample period and a citation for the federal Export Diversification Strategy
(2018) and subsequent commitments.}}
Yet countries do not choose their trading partners; firms do. A nation's trade
portfolio is the sum of many firm-level decisions about which foreign
suppliers to buy from, which foreign customers to sell to, and in which
products. Understanding what moves a firm into a new market is therefore a
prerequisite for understanding how national trade patterns change.

Entering a foreign market requires knowledge that is costly to acquire: who
the reliable suppliers and buyers are, how business is conducted, how
contracts are written and enforced, and how to navigate an unfamiliar
language and set of institutions. Much of this knowledge is held by people
rather than by organizations. Immigrants are a natural source of it. In 2021,
23 percent of Canada's population was born abroad, and a large literature
shows that countries trade more with the places their immigrants come from
\citep{gould1994,headries1998,rauchtrindade2002}. Recent firm-level work
locates part of this pattern inside businesses: Canadian firms with immigrant
owners import more from, and export more to, their owners' countries of
origin \citep{fung2019,fungliu2025}, and U.S. firms founded by immigrants
source disproportionately from their founders' birth countries, where they pay
lower prices \citep{chikissaygin2026}.

These correlations admit two readings with very different implications. On
the first, owners carry knowledge that lowers their firm's cost of trading
with their country of origin; who owns a firm then shapes what it trades, and
the owners a country's businesses attract become a lever on its trade
portfolio. On the second, immigrants found or acquire firms that are already
well placed to trade with their origin countries---because of the products
the firm makes, the customers it serves, or the networks it already
uses---and ownership merely labels firms that would have traded with those
countries anyway. Cross-sectional comparisons cannot separate the two. And
even under the first reading, it matters \emph{which} cost the owner lowers.
If the owner makes it cheaper to start new relationships, the firm diversifies
into new partners and products. If the owner instead makes existing
relationships cheaper to run, trade deepens along lines the firm had already
established. If the owner chiefly shifts demand toward home-country goods,
imports are reoriented without any change in the cost of trading.

This paper asks whether a change in ownership from a non-immigrant to an
immigrant owner changes an incumbent firm's international trade, and through
which of these channels. I study roughly 1,000 Canadian firms whose ownership
passes from non-immigrant to immigrant owners between 2005 and 2019. I link
Statistics Canada's firm-level customs records---imports by six-digit
Harmonized System (HS6) product and source country since 2001, exports by
product and destination since 2005---to owner records that identify immigrant
owners and their countries of origin. The key feature of the setting is that
the firm, its products, its location, and its existing relationships are held
fixed while the person in control changes. Any change in where the firm trades
that coincides with the arrival of an owner from country $o$, and that is
specific to country $o$, is difficult to attribute to the business rather
than to the person.

The empirical design compares, within the same firm and year, the change in
trade with the new owner's country of origin to the change in trade with all
other countries, relative to the same comparison in firms that did not change
hands or were acquired by owners from other origins. Firm-by-year fixed
effects absorb everything that changes when a firm is sold---new management,
new financing, reorganization, changes in scale---provided it does not
differentially favour country $o$. Country-by-product-by-year fixed effects
absorb shocks to Canada's trade with $o$ in the firm's products, including
tariff changes under trade agreements signed during the period and exchange
rate movements. Firm-by-country-by-product fixed effects absorb persistent
differences in a firm's fit with particular partners, including any selection
of buyers on the \emph{level} of a firm's opportunities with their origin
country. The remaining threat is that buyers select firms in anticipation of
firm- and country-specific shocks that would have occurred without them. I
assess it with event-study leads, comparisons among firms acquired by owners
from different origins, comparisons with transfers between non-immigrant
owners, placebo origin countries, and the timing of sales initiated by
sellers.

To interpret the estimates, I develop a simple partial-equilibrium model in
which a firm chooses which country--product trade relationships to form, in
each direction, and how much to trade through each. Forming a relationship
requires a fixed cost at the product level and, the first time the firm trades
with a country, a fixed cost at the country level; trading through it incurs a
variable cost. Each friction has a component that every Canadian firm faces
and an information component that an owner from the partner country can
partly remove. A parallel channel allows the owner to raise demand for
home-country consumer goods. The model maps the three channels into distinct
predictions: a fixed-cost channel generates entry into new country--product
pairs, in bundles when the country-level cost falls, and leaves incumbent
flows unchanged; a variable-cost channel raises incumbent flows in proportion;
and a demand channel affects only imports of consumer goods, leaving exports
and imported intermediate inputs untouched. My leading hypothesis is the
fixed-cost channel. Information about \emph{whom} to trade with is acquired
once, at the start of a relationship; search frictions in international goods
markets are large \citep{lenoir2023,startz2025}; and the U.S. evidence that
immigrant founders reach different, lower-priced suppliers rather than
better terms from the same suppliers \citep{chikissaygin2026} points to an
advantage in finding partners. The design and the model, however, treat the
three channels symmetrically, and the data decide among them.

The model delivers two further results that shape the empirical analysis.
First, when sourcing decisions are interdependent, as in
\citet{antras2017}, an owner's advantage in one country changes the firm's
scale and therefore its trade with every other country. Under the
interdependence structure of that model, these spillovers are common across
the firm's other partners, so the within-firm comparison continues to
identify the country-specific effect, while trade with third countries
requires a separate comparison across firms. Second, the model clarifies what
firm-level diversification means: an expansion of trade with the owner's
origin raises the number of the firm's partners and relationships, but lowers
the concentration of its trade across partners only when the origin country's
initial share is below the firm's Herfindahl index. Because most Canadian
traders are heavily concentrated on the United States, the condition is
likely to hold for most firms and origins; whether it does, and whether the
new trade is created or diverted from other partners, is an empirical
question that links the firm-level estimates to national diversification.

\result{Paragraph 1 of findings: event-study evidence on trade with the
origin country for imports and exports; size of the pooled effect; absence
or presence of pre-trends.}

\result{Paragraph 2 of findings: decomposition into country entry, new
country--product pairs, and incumbent flows; which channel the evidence
favours; results for firms already trading with the origin country before the
switch; consumer versus intermediate goods and imports versus exports.}

\result{Paragraph 3 of findings: firm-level diversification (number of
partners, concentration, non-U.S. share); creation versus diversion with
respect to third countries; robustness and placebo tests.}

This paper contributes to three literatures. The first studies how migrant
networks shape trade. Since \citet{gould1994} and \citet{headries1998}, who
used Canadian data, a large body of work has related bilateral trade to
immigrant stocks across countries and regions
\citep{wagner2002,rauchtrindade2002,combes2005,peri2010}, and quasi-experimental
studies using refugee settlement and other plausibly exogenous variation in
migrant populations have confirmed a causal link
\citep{cohen2017,parsonsvezina2018,steingress2018,bonadio2026}. Because these
studies measure trade and migration at the level of countries or regions, they
cannot say which firms carry the trade, through which decisions, or which
friction migrants relax---whether information costs or preferences, a
distinction emphasized since \citet{rauch2001}. I move the analysis inside the
firm and use a change in the person who controls it, rather than a change in
the local population, to identify the effect.

The second literature studies immigrants inside firms. A series of papers
shows that firms employing immigrants trade more with immigrants' origin
countries
\citep{hiller2013,hatzigeorgiou2016,andrews2017,mitaritonna2017,ottaviano2018,egger2019,ariu2022,erbahar2023},
including Canadian evidence on exports \citep{cardoso2022}. Closest to this
paper, \citet{fung2019} and \citet{fungliu2025} document that Canadian firms
with immigrant owners trade more with their owners' origins, and
\citet{chikissaygin2026} show that U.S. firms founded by immigrants source
from their founders' birth countries at lower prices and quantify the value of
that advantage in a model of supplier search and entrepreneurship. These
papers compare firms with and without immigrant owners or founders. A central
assumption of the quantitative work---that the trade advantage is embodied in
the person, so that it would disappear if the founder faced native costs---is
not tested by cross-sectional variation, because the founder and the business
she designs are observed together. Holding the business fixed while the owner
changes provides a direct test of portability, and it does so for incumbent
firms, the population that accounts for most trade but that founder-based
designs observe least. The paper also relates to work showing that
individuals carry market-specific knowledge across firms: managers moving
between firms \citep{mionopromolla2014}, sales managers transmitting
foreign customer capital \citep{patault2024}, employee spinoffs entering their
parents' export markets \citep{muendler2018}, and firms learning to import from
peers \citep{bisztray2018}. More broadly, it relates to evidence that the
identity of those who control a firm affects its decisions and performance
\citep{bertrandschoar2003,bennedsen2007,becker2022}, and that changes in
ownership change firm behaviour \citep{guadalupe2012}. Here the organizational
event is a change in ownership and the outcome is the geography of trade.

The third literature studies the margins of firm-level trade and the frictions
behind them. Fixed and sunk costs of market entry shape which firms trade
and where \citep{melitz2003,chaney2008,das2007,eaton2011}, firms expand
sequentially across products and markets
\citep{bernard2007,bernard2011,albornoz2012,morales2019}, and sourcing decisions
are interdependent across countries \citep{antras2017}. A growing literature
attributes part of these costs to information and search frictions
\citep{allen2014,chaney2014,steinwender2018,lenoir2023,startz2025} and shows
that trade relationships, once formed, are valuable and persistent
\citep{monarch2023}. Distinguishing fixed from variable trade costs is
difficult because both are typically inferred from the same flows. An owner's
country of origin provides an observed, person-specific, and country-specific
shifter of bilateral frictions within a firm, and the timing of ownership
changes allows me to trace how a firm's relationship portfolio responds. By
linking these firm-level responses to concentration and partner composition,
the paper also connects to work on the diversification of national trade
\citep{imbs2003,koren2007,caselli2020}.

The remainder of the paper proceeds as follows. Section~\ref{sec:data}
describes the setting and data. Section~\ref{sec:model} presents the model and
its predictions. Section~\ref{sec:strategy} describes the empirical strategy.
Sections~\ref{sec:results} and~\ref{sec:mechanisms} present the main results
and the mechanism tests. Section~\ref{sec:discussion} discusses
interpretation, aggregation, and policy, and Section~\ref{sec:conclusion}
concludes.

%=====================================================================
\section{Setting and Data}
\label{sec:data}
%=====================================================================

\subsection{Setting}

Canada is a useful setting for three reasons. First, immigrants are a large
share of the population and of business owners, and they come from a wide
range of origin countries, which generates variation in the origin of
incoming owners across firms. \verify{Share of private enterprises with
immigrant majority ownership; citation to Statistics Canada.} Second, the
country's trade is highly concentrated on one partner, so that firm-level
moves into other markets are directly relevant to policy. Third, a large
number of privately held businesses change hands each year, many as owners
retire, which creates the ownership transitions the design requires.
\verify{Citation on business succession and retirement-driven sales.}

During the sample period Canada implemented several trade agreements that
changed tariffs with particular partners, including the Canada--Korea Free
Trade Agreement (2015), the provisional application of the Comprehensive
Economic and Trade Agreement with the European Union (2017), and the
Comprehensive and Progressive Agreement for Trans-Pacific Partnership (2018).
Such changes affect all Canadian firms trading a given product with a given
country and are absorbed by country-by-product-by-year fixed effects in the
main specification.

\subsection{Trade data}

The trade data are Statistics Canada's Trade by Importer Characteristics (TIC)
and Trade by Exporter Characteristics (TEC) microdata, which link customs
transactions to business identifiers. Imports are observed from 2001 to 2019
by HS6 product and country of origin; exports are observed from 2005 to 2019
by HS8 product and country of destination. I aggregate both to the
firm--country--HS6--year level, harmonizing export codes to HS6 and concording
HS revisions (2002, 2007, 2012, 2017) to a time-consistent classification. The
data record the value of trade in each cell; physical quantities are not used.
\verify{Whether imports are recorded by country of origin or country of
shipment; whether domestic exports and re-exports can be separated; reporting
thresholds; treatment of low-value and courier shipments.}

The main analysis uses 2005--2019 for both directions so that imports and
exports cover the same window. Imports for 2001--2004 extend the pre-period
for early switchers and are used in robustness checks. Ending in 2019 avoids
the disruption to trade during the COVID-19 pandemic. Values are deflated to
constant dollars. \verify{Deflator.}

\subsection{Ownership and immigrant status}

Ownership information comes from owner-linked business records, following the
construction in \citet{fung2019}. \verify{Source of ownership records (e.g.,
shareholder schedules for incorporated firms and individual tax records for
unincorporated businesses), threshold for recording an owner, and linkage to
immigrant landing records.} An owner is classified as an immigrant if she
appears in the immigrant landing records, and her country of origin is
\verify{country of birth / citizenship / last permanent residence}. A firm is
classified as immigrant-owned in a year if \verify{immigrant owners together
hold a majority of the firm's shares; alternatively, if the largest
shareholder is an immigrant}.

Two measurement issues matter for the design. First, immigrants who landed
before the start of the landing records are classified as non-immigrants.
\verify{Start year of landing-record coverage (believed to be 1980).} A firm
sold by an early-cohort immigrant to a recent immigrant from the same country
would then appear as a non-immigrant-to-immigrant transition. I therefore
describe the baseline group as non-immigrant ``as measured'' and test
sensitivity to excluding sellers old enough to have landed before coverage
began. Second, firms with several owners may have owners from more than one
origin. The baseline assigns the origin of the owner or owners holding the
largest combined share; firms whose incoming owners come from several origins
are analysed separately.

\subsection{Firm continuity and the event definition}

A change in ownership must be distinguished from administrative events that
mechanically alter reported trade. I require that the firm's business
identifier persist across the transition and exclude transitions that coincide
with mergers, amalgamations, consolidations of reporting units, or changes in
the identifier under which customs declarations are filed. \verify{Available
indicators of mergers and identifier changes.}

The event is the first year $t_i^*$ in which firm $i$ moves from non-immigrant
to immigrant ownership. The incoming owner's origin $o_i$ is fixed throughout
the event window, including the pre-period. The main sample retains firms that
traded in at least one direction in the pre-period and that are observed in
the business register for at least \verify{two} years before the transition.
Firms are \emph{not} required to survive or to trade after the transition:
conditioning on post-transition outcomes would select on a possible effect of
treatment. Firms that exit after the transition contribute zeros where
business-register coverage indicates they are no longer operating, and are
otherwise treated as missing. Repeated transitions and reversals are flagged;
in the baseline, firms are followed until a second change of ownership.

Table~\ref{tab:switches} will report the distribution of qualifying
transitions by year, origin region, industry, and pre-period trading status.
The sample contains roughly 1,000 qualifying transitions. \verify{Exact count
after all restrictions; disclosure rounding.}

\begin{table}[htbp]\centering\small
\caption{Qualifying ownership transitions (shell)}
\label{tab:switches}
\begin{tabular}{lccc}
\toprule
 & Transitions & Share importing before & Share exporting before\\
\midrule
\multicolumn{4}{l}{\textit{Panel A. By origin region of incoming owner}}\\
\quad East Asia & -- & -- & -- \\
\quad South Asia & -- & -- & -- \\
\quad Europe & -- & -- & -- \\
\quad Middle East and Africa & -- & -- & -- \\
\quad Americas (excl.\ U.S.) & -- & -- & -- \\
\quad United States & -- & -- & -- \\
\multicolumn{4}{l}{\textit{Panel B. By sector}}\\
\quad Manufacturing & -- & -- & -- \\
\quad Wholesale trade & -- & -- & -- \\
\quad Retail trade & -- & -- & -- \\
\quad Other & -- & -- & -- \\
\midrule
Total & $\approx$1,000 & -- & -- \\
\bottomrule
\end{tabular}
\end{table}

\subsection{Comparison groups}

Three groups of firms serve as comparisons: (i) firms that remain under
non-immigrant ownership throughout the window; (ii) firms that pass from one
non-immigrant owner to another, which share the disruption of a sale but not
the arrival of an immigrant owner; and (iii) firms that pass to immigrant
owners from a different origin, which share the arrival of an immigrant owner
but not the specific origin. Comparisons within groups (ii) and (iii) are the
most demanding, because they hold fixed the selection of firms into being sold
and, in (iii), the selection of firms into immigrant ownership.

\subsection{Outcomes and the panel}

For each direction $d \in \{\text{import}, \text{export}\}$, I construct a
panel at the firm--country--product--year level that retains zeros for
relationships the firm could plausibly have traded. Because the full set of
firm--country--product cells is very large, the set of potential relationships
for firm $i$ is defined from pre-period information only: all HS6 products the
firm traded with any country in the pre-period, together with the products
most commonly traded by firms in its six-digit industry, crossed with all
partner countries. Absence from the customs records is coded as zero only
where business-register coverage indicates that the firm was operating.

The outcomes are:
\begin{itemize}[leftmargin=1.5em,itemsep=0.1em]
\item \emph{Bilateral value}: trade with country $c$ in product $k$, $X_{ickdt}$,
and its sum over products, $X_{icdt}$.
\item \emph{Country participation}: $\one[X_{icdt}>0]$.
\item \emph{Product breadth}: the number of products traded with country $c$.
\item \emph{New and incumbent relationships}: trade in country--product pairs
not active in the pre-period, and trade in pairs active in the year before the
transition (with zeros retained after the transition).
\item \emph{Firm-level diversification}: the number of partner countries and of
country--product relationships; the Herfindahl index of trade across partner
countries, $H_{idt} = \sum_c s_{icdt}^2$; and the non-U.S.\ share of trade.
\item \emph{Totals}: total imports and exports, and trade with countries other
than the owner's origin.
\end{itemize}

Products are classified as consumer, intermediate, or capital goods using the
UN Broad Economic Categories, and as differentiated, reference-priced, or
homogeneous using \citet{rauch1999}. Country characteristics---distance,
common official language with Canada, institutional quality
\citep{kaufmann2011}, and the size of Canada's aggregate trade with the
country---are used to measure the information content of bilateral frictions.

\subsection{Descriptive statistics}

\begin{table}[htbp]\centering\footnotesize
\caption{Pre-period characteristics of switching and comparison firms (shell)}
\label{tab:sumstats}
\begin{tabular}{lcccc}
\toprule
 & Treated & Sold to & Never & Sold to imm.\\
 & & non-imm. & sold & (other origin)\\
\midrule
Employment & -- & -- & -- & -- \\
Importer (share) & -- & -- & -- & -- \\
Exporter (share) & -- & -- & -- & -- \\
Import partner countries & -- & -- & -- & -- \\
Export partner countries & -- & -- & -- & -- \\
HS6 products imported & -- & -- & -- & -- \\
HS6 products exported & -- & -- & -- & -- \\
Trades with $o$ (share) & -- & \multicolumn{3}{c}{-- (pseudo-origin)} \\
Share of trade with $o$ & -- & \multicolumn{3}{c}{-- (pseudo-origin)} \\
HHI across partners & -- & -- & -- & -- \\
Non-U.S.\ share of trade & -- & -- & -- & -- \\
\midrule
Firms & $\approx$1,000 & -- & -- & -- \\
\bottomrule
\end{tabular}
\end{table}

\planned{Planned descriptive facts: (i) the cross-sectional origin gradient
among all firms, replicating \citet{fungliu2025} for both directions and
splitting it by margin; (ii) how switching firms compare with other firms
before the transition; (iii) how often switching firms already trade with the
incoming owner's origin before the transition.}

%=====================================================================
\section{A Model of Trade Relationships and Owner Knowledge}
\label{sec:model}
%=====================================================================

The model is deliberately simple. Its purpose is to state which cost an owner
could change, and what each possibility implies for each margin of trade,
rather than to provide a new general theory of firm-level trade. It takes the
firm's operations in its industry as given and models its choice of trade
relationships.

\subsection{Environment}

Firm $i$ has an owner with origin $o_i$; for a non-immigrant owner,
$o_i=\varnothing$. A trade \emph{relationship} is a triple $(c,k,d)$: a
partner country $c$, an HS6 product $k$, and a direction
$d\in\{\text{im},\text{ex}\}$. An active relationship earns variable profit
and pays a fixed cost:
\begin{equation}
\pi_{ickd} = \Phi_{ickd}\,\tau_{icd}^{-\varepsilon_d} - f_{ickd}.
\label{eq:pi}
\end{equation}
$\Phi_{ickd}$ collects the firm's scale and productivity, supply or demand
conditions for product $k$ in country $c$, and prices. $\tau_{icd}\geq1$ is the
firm's variable cost of trading with $c$---transport, contracting, monitoring,
and enforcement. $f_{ickd}$ is the fixed cost of operating the
relationship---finding a partner, learning product-specific procedures,
writing a first contract. $\varepsilon_d>0$ is the elasticity of variable profit
with respect to the variable cost. In addition, the firm pays a country-level
fixed cost $F_{icd}$ in any period in which it has at least one active
relationship with $c$ in direction $d$. It covers learning the country's
institutions, language, payment practices, and customs procedures.

The firm chooses, for each direction, the set of active countries
$\mathcal{C}_{id}$ and, for each active country, the set of active products
$\mathcal{K}_{icd}$, to maximize
\begin{equation}
\Pi_i = \sum_{d}\sum_{c\in\mathcal{C}_{id}}\Big[\sum_{k\in\mathcal{K}_{icd}}
\pi_{ickd} - F_{icd}\Big].
\label{eq:Pi}
\end{equation}

\begin{assumption}[Separability and trade values]\label{ass:sep}
(i) Relationships are separable: $\Phi_{ickd}$ does not depend on the firm's
other relationships. (ii) The value of trade through an active relationship is
proportional to its variable profit, $X_{ickd}=\mu_d\,\Phi_{ickd}\,
\tau_{icd}^{-\varepsilon_d}$, with $\mu_d>0$ common across firms, countries,
and products.
\end{assumption}

Part (ii) holds in each of the sectors the data cover. For an exporter facing
constant-elasticity demand abroad, revenue is proportional to variable
profit, with $\varepsilon=\sigma-1$. For a wholesaler or retailer importing
goods for resale under constant-elasticity demand at home, the value of
purchases is proportional to the variable profit from reselling them. For a
manufacturer importing inputs, input spending is proportional to the
variable profit the input generates when sourcing decisions are separable;
Section~\ref{sec:interdep} relaxes separability for this case. The firm's
type therefore affects only $\Phi$ and $\varepsilon$, not the mechanism.

\begin{lemma}[Optimal relationship set]\label{lem:set}
Under Assumption~\ref{ass:sep}(i), conditional on country $c$ being active in
direction $d$, product $k$ is active if and only if
$\Phi_{ickd}\tau_{icd}^{-\varepsilon_d}\geq f_{ickd}$. Country $c$ is active
if and only if
\begin{equation}
\Omega_{icd}\equiv\sum_{k}\max\{\Phi_{ickd}\tau_{icd}^{-\varepsilon_d}-f_{ickd},\,0\}
\;\geq\;F_{icd}.
\label{eq:Omega}
\end{equation}
\end{lemma}

\noindent The proof is immediate from \eqref{eq:Pi}: the problem separates
across directions and countries, and within a country across products.

\subsection{What an immigrant owner changes}

Let $D_{ic}=\one[c=o_i]$. Each friction has a component every Canadian firm
faces and an \emph{information component} that an owner from the partner
country can partly remove. An owner from $o$ removes a share of the
information component for her country of origin only:
\begin{align}
\ln\tau_{icd} &= \ln\bar\tau_{cd} + (1-\lambda_\tau D_{ic})\,\iota_{c},
\label{eq:tau}\\
F_{icd} &= \bar F_{cd} + (1-\lambda_F D_{ic})\,\psi_{c},
\label{eq:F}\\
f_{ickd} &= \bar f_{ckd} + (1-\lambda_f D_{ic})\,\psi_{ck},
\label{eq:f}
\end{align}
with $\lambda_\tau,\lambda_F,\lambda_f\in[0,1]$ and information components
$\iota_c,\psi_c,\psi_{ck}\geq0$.

\begin{assumption}[Information components]\label{ass:info}
The information components $\iota_c$, $\psi_c$, and $\psi_{ck}$ are larger where
public information about trading with $c$ is scarce: for differentiated
goods, for countries with weaker institutions or greater linguistic and
institutional distance from Canada, and for countries with which Canada's
aggregate trade is thin.
\end{assumption}

Assumption~\ref{ass:info} is the source of testable heterogeneity. If the owner
instead lowered $\tau$ by a constant proportion everywhere, the proportional
effect would be the same for every country and product, and there would be no
cross-sectional pattern to test.

\paragraph{Demand channel.} For imports, the owner may also shift demand
toward home-country goods, for example by serving co-ethnic customers:
\begin{equation}
\Phi_{ick,\text{im}} = \bar\Phi_{ick,\text{im}}\,\big(1+\eta\,D_{ic}\,\one[k\in\mathcal{K}^{\text{cons}}]\big),
\label{eq:demand}
\end{equation}
where $\mathcal{K}^{\text{cons}}$ is the set of consumer goods and $\eta\geq0$.

\begin{assumption}[Scope of the demand channel]\label{ass:demand}
The demand channel affects imports of consumer goods only. It does not affect
exports to $o$ or imports of intermediate and capital goods.
\end{assumption}

Assumption~\ref{ass:demand} is what makes the demand channel testable, and it
is also its weakest point. Two departures are plausible: firms that sell
food or other inputs to co-ethnic businesses may import intermediate goods
for that demand, and an owner may prefer home-country suppliers for reasons of
trust rather than cost. Both would load onto the cost channels in the tests
below. I return to these possibilities in Section~\ref{sec:mechanisms}.

\subsection{The effect of a change in ownership}

Consider firm $i$ whose owner changes at $t^*$ from a non-immigrant to an
owner from $o$, holding all primitives other than $D_{io}$ fixed. Write
$\Delta$ for the change across the transition. The following results describe
each channel in isolation; when several operate, their effects combine.

\begin{proposition}[Bilateral effect]\label{prop:bilateral}
Under Assumption~\ref{ass:sep}, trade with $o$ weakly increases under each
channel, and trade with every $c\neq o$ is unchanged. The within-firm change in
trade with $o$ relative to other countries therefore equals the change in
trade with $o$.
\end{proposition}

\begin{proposition}[Margins]\label{prop:margins}
Under Assumption~\ref{ass:sep}:
\begin{enumerate}[label=(\alph*),itemsep=0.1em]
\item \emph{Variable cost} ($\lambda_\tau>0$ only). Every incumbent relationship
with $o$ grows by the same proportion,
$\Delta\ln X_{iokd}=\varepsilon_d\lambda_\tau\iota_o$. New relationships enter
only if they were within that proportion of the activation threshold.
\item \emph{Product fixed cost} ($\lambda_f>0$ only). Incumbent relationships are
unchanged. Trade with $o$ grows only through new pairs $(o,k)$ with
$\bar f_{okd}+(1-\lambda_f)\psi_{ok}\leq\Phi_{iokd}\tau_{iod}^{-\varepsilon_d}
<\bar f_{okd}+\psi_{ok}$, provided $o$ is or becomes active.
\item \emph{Country fixed cost} ($\lambda_F>0$ only). Firms already active with
$o$ before the transition are unaffected. A firm not previously active with $o$
enters if and only if $\bar F_{od}+(1-\lambda_F)\psi_o\leq\Omega_{iod}<
\bar F_{od}+\psi_o$; when it does, it activates at once every product with
positive net profit with $o$.
\item \emph{Demand} ($\eta>0$ only). Incumbent consumer-good import
relationships with $o$ grow by $\ln(1+\eta)$; new consumer-good import
relationships may enter. Exports to $o$ and imports of intermediate and
capital goods from $o$ are unchanged.
\end{enumerate}
\end{proposition}

\noindent The proof, in Appendix~\ref{app:proofs}, follows from
Lemma~\ref{lem:set} and Assumption~\ref{ass:sep}(ii). Proposition~\ref{prop:margins}
is the core of the mechanism tests. Under a fixed-cost channel the response is
entirely extensive, and under the country-level component it is concentrated
among firms new to $o$ and arrives in bundles of products. Under a
variable-cost channel incumbent flows grow in proportion. Under a demand
channel the response is confined to imports of consumer goods.

\begin{proposition}[Heterogeneity]\label{prop:hetero}
Under Assumptions~\ref{ass:sep}--\ref{ass:info}, the proportional response of
incumbent relationships under the variable-cost channel,
$\varepsilon_d\lambda_\tau\iota_o$, is increasing in $\iota_o$. Under the
fixed-cost channels, the set of relationships pushed above the activation
threshold widens with $\psi_{ok}$ and $\psi_o$. Under the demand channel,
responses are unrelated to $\iota$ and $\psi$.
\end{proposition}

\noindent A wider window does not guarantee more entry: the number of
newly active relationships also depends on how many relationships lie near
the threshold. Where frictions are very high, firms may remain too far below
the threshold for the owner's knowledge to make trade profitable. The extensive
response to $\psi$ can therefore be non-monotone, and the empirical tests
treat it as a prediction about the incumbent margin first and the extensive
margin second.

\subsection{Firm-level diversification}

The model also clarifies what an owner's arrival does to the diversification
of a firm's trade.

\begin{proposition}[Diversification]\label{prop:divers}
Under Assumption~\ref{ass:sep}:
\begin{enumerate}[label=(\roman*),itemsep=0.1em]
\item The number of the firm's partner countries and of its country--product
relationships weakly increases.
\item Let $s_{io}$ be the share of country $o$ in the firm's trade in direction
$d$ and $H_i=\sum_c s_{ic}^2$ the Herfindahl index across partner countries. A
marginal increase in trade with $o$, holding trade with other countries fixed,
lowers $H_i$ if and only if $s_{io}<H_i$.
\item If $o$ is not the United States, the firm's non-U.S.\ share of trade
weakly increases.
\end{enumerate}
\end{proposition}

\noindent Part (ii) is the relevant notion of diversification for a firm that
already trades. An expansion toward a partner that is small relative to the
firm's existing concentration diversifies the firm; an expansion that makes
$o$ the dominant partner concentrates it. For a firm whose trade is dominated
by the United States, $H_i$ is close to the U.S.\ share, and almost any
expansion toward a non-U.S.\ origin diversifies. Proposition~\ref{prop:divers}
relies on separability. If trade with $o$ comes at the expense of other
partners, the effect on diversification depends on which partners lose---the
subject of the next subsection.

\subsection{Interdependence across countries}
\label{sec:interdep}

Separability is restrictive for manufacturers importing inputs, because adding
a lower-cost source lowers the firm's marginal cost and raises its scale. To
capture this, I adopt the sourcing structure of \citet{antras2017}. Firm $i$
with productivity $\varphi_i$ sources inputs from a set of countries
$\mathcal{J}_i$ and has sourcing capability
\begin{equation}
\Theta_i=\sum_{j\in\mathcal{J}_i}T_j\,(\tau_{ij}w_j)^{-\theta},
\label{eq:Theta}
\end{equation}
where $T_j$ is country $j$'s state of technology, $w_j$ its input cost, and
$\theta$ the dispersion of input efficiency across countries. Its unit cost is
proportional to $\varphi_i^{-1}\Theta_i^{-1/\theta}$, its share of input
purchases from $j$ is $T_j(\tau_{ij}w_j)^{-\theta}/\Theta_i$, and its total
input purchases are proportional to $\Theta_i^{(\sigma-1)/\theta}$, where
$\sigma$ is the elasticity of demand for the firm's output.

\begin{proposition}[Interdependence]\label{prop:interdep}
In the sourcing model of \eqref{eq:Theta}, holding the sourcing set fixed,
imports from country $j$ satisfy
$X_{ij}\propto T_j(\tau_{ij}w_j)^{-\theta}\,\Theta_i^{(\sigma-1)/\theta-1}$.
An owner who lowers $\ln\tau_{io}$ by $\lambda_\tau\iota_o$:
\begin{enumerate}[label=(\alph*),itemsep=0.1em]
\item raises imports from every $j\neq o$ in the same proportion if
$(\sigma-1)/\theta>1$ and lowers them in the same proportion if
$(\sigma-1)/\theta<1$;
\item changes imports from $o$ relative to any $j\neq o$ by exactly
$\theta\lambda_\tau\iota_o$ in logs.
\end{enumerate}
When $(\sigma-1)/\theta>1$, sourcing decisions are complements and the set of
source countries weakly expands \citep{antras2017}.
\end{proposition}

\noindent Proposition~\ref{prop:interdep} has two implications for the
empirical design. First, the effect on third countries is common to all of
the firm's other partners. A within-firm comparison with firm-by-year fixed
effects absorbs it and identifies the country-specific effect in part (b)
without requiring separability. Second, because the third-country effect is
absorbed, measuring it---and hence distinguishing trade creation from trade
diversion---requires a comparison across firms. Whether the owner's
arrival diversifies the firm then depends on whether third-country trade
rises (complements) or falls (substitutes), and on which partners lose.
For exporters, interdependence may arise through capacity constraints, which
make destinations substitutes; the same logic applies.

\subsection{Dynamics and persistence}

Two extensions sharpen the timing predictions. First, if fixed costs are
sunk---paid once on entry, with lower costs of maintaining a relationship
\citep{das2007}---and finding a partner takes time, the fixed-cost channel
builds over several years after the transition, while a variable-cost
reduction acts immediately. Second, under sunk costs, relationships formed
under an immigrant owner persist after she leaves, while the gains from a
lower variable cost leave with her. With roughly 1,000 transitions, reversals
are too rare to support a main-text test, and the persistence prediction is
treated as exploratory.

\subsection{Selection into ownership}
\label{sec:selection}

Ownership is not randomly assigned. A buyer from $o$ values firm $i$ at its
baseline value $V_i$ plus the gain from her knowledge,
\begin{equation}
G_i(o)=\sum_{t\geq t^*}\beta^{t-t^*}\big[\Pi_{it}(D_{io}=1)-\Pi_{it}(D_{io}=0)\big],
\end{equation}
which is larger for firms whose products $o$ supplies or demands strongly
(high $\Phi_{iok}$) and where information components are large. Buyers from
$o$ therefore select firms with high $G_i(o)$.

\begin{proposition}[What selection does and does not threaten]\label{prop:selection}
(i) Because buyers select on $G_i(o)$, the estimand is the average effect of
immigrant ownership on the firms immigrants acquire (the ATT), not the average
effect across all firms. (ii) Selection on the time-invariant components of
$\Phi_{iokd}$ and of the frictions is absorbed by firm-by-country-by-product
fixed effects. (iii) The design fails if buyers select on anticipated changes
in $\Phi_{iokd,t}$---shocks to the firm's opportunities with $o$ that would have
occurred without the change in ownership.
\end{proposition}

\noindent Part (iii) is the identifying assumption stated in economic terms.
It is violated if, for example, a buyer acquires a firm because a contract
with a customer in $o$ is already in prospect. It is not violated if a buyer
acquires a firm because its products are ones $o$ supplies cheaply, which is
selection on the level of the gain. A related form of selection is through
networks: a firm already trading with $o$ is more likely to come to the
attention of buyers from $o$. This affects which firms are treated, and so the
interpretation of heterogeneity by pre-transition trading status, but not
the validity of the main comparison.

\subsection{Summary of predictions}

Table~\ref{tab:predictions} summarizes the predictions of each channel and the
test that corresponds to each.

\begin{table}[htbp]\centering\footnotesize
\caption{Predictions by channel}
\label{tab:predictions}
\renewcommand{\arraystretch}{1.25}
\begin{tabularx}{\textwidth}{>{\raggedright\arraybackslash}p{3.3cm}YYYY}
\toprule
Prediction & Country fixed cost ($\lambda_F$) & Product fixed cost ($\lambda_f$) & Variable cost ($\lambda_\tau$) & Demand ($\eta$)\\
\midrule
P1. Trade with $o$ relative to other countries & Rises & Rises & Rises & Rises (consumer imports only)\\
P2. Incumbent $(o,k)$ flows & Unchanged & Unchanged & Rise by $\varepsilon\lambda_\tau\iota_o$ & Consumer imports rise\\
P3. New $(o,k)$ pairs & Firms new to $o$ only & Yes & At threshold only & Consumer imports only\\
P4. Firms already trading with $o$ & No effect & Effect & Effect & Effect\\
P5. Breadth of entry & Many products at once & One at a time & Few & Consumer goods only\\
P6. Exports and intermediate imports & Rise & Rise & Rise & No effect\\
P7. Gradient in information intensity & Larger where $\psi$ large & Larger where $\psi$ large & Larger where $\iota$ large & None\\
P8. Timing & Gradual (search) & Gradual (search) & Immediate & Immediate\\
\bottomrule
\end{tabularx}
\end{table}

%=====================================================================
\section{Empirical Strategy}
\label{sec:strategy}
%=====================================================================

\subsection{Estimand and identifying assumption}

Let $X_{ickdt}(1)$ and $X_{ickdt}(0)$ denote potential trade with and without a
transition to an owner from $o_i$. The target parameter is the proportional
effect of the transition on trade with $o_i$ among switching firms, relative to
the effect on the same firms' trade with other countries. Because trade
values contain many zeros and the model is multiplicative, I state the
identifying assumption in ratios \citep{wooldridge2023}: absent the
transition, the ratio of a switching firm's expected trade with $o_i$ to its
expected trade with other countries would have evolved as in comparison firms,
conditional on the fixed effects below. In the language of
Proposition~\ref{prop:selection}, buyers may select firms on the level of
their opportunities with $o$ but not on anticipated changes in them.

\subsection{Main specification}

For each direction $d$, I estimate by Poisson pseudo-maximum likelihood
\citep{santossilva2006}, retaining zeros,
\begin{equation}
X_{ickt} = \exp\big(\beta\,\Post_{it}\times\one[c=o_i] + \alpha_{it} +
\gamma_{ckt} + \delta_{ick}\big)\,\epsilon_{ickt},
\label{eq:main}
\end{equation}
where $\Post_{it}=\one[t\geq t_i^*]$ for switching firms. The coefficient
$\beta$ is the proportional change in trade with the owner's origin caused by
the transition, relative to other countries; $\exp(\beta)-1$ is the
percentage effect. Each fixed effect has a counterpart in the model
(Table~\ref{tab:fe}). PPML is consistent under the multiplicative model,
retains the zeros that the extensive margin requires, and weights
observations by their expected values; I report results from unweighted
linear models in levels and participation alongside it. High-dimensional
fixed effects are estimated following \citet{correia2020}.

\begin{table}[htbp]\centering\small
\caption{Fixed effects in equation~\eqref{eq:main} and their model counterparts}
\label{tab:fe}
\renewcommand{\arraystretch}{1.2}
\begin{tabularx}{\textwidth}{>{\raggedright\arraybackslash}p{3.6cm}YY}
\toprule
Fixed effect & Absorbs & Model counterpart\\
\midrule
Firm $\times$ year ($\alpha_{it}$) & Any general effect of the change in ownership: new management, finance, reorganization, firm-level demand and scale & Firm component of $\Phi$; common third-country effects under interdependence (Prop.~\ref{prop:interdep})\\
Country $\times$ product $\times$ year ($\gamma_{ckt}$) & Supply and demand shocks in $c$, exchange rates, tariffs and trade agreements, gravity & Partner component of $\Phi$; common component of $\tau$\\
Firm $\times$ country $\times$ product ($\delta_{ick}$) & Persistent firm--partner fit, including selection of buyers on the level of opportunities with $o$ & Time-invariant components of $\Phi$ and of frictions (Prop.~\ref{prop:selection}(ii))\\
\bottomrule
\end{tabularx}
\end{table}

\subsection{Event study}

To trace the time path and assess pre-trends, I replace $\Post_{it}$ with
indicators for event time $e=t-t_i^*$, binned at the endpoints of a window
$e\in[-4,6]$ \verify{window chosen given 2005--2019 coverage and transition
years}. Because transitions are staggered and effects may vary across cohorts
and over time, two-way fixed-effects estimates can be biased
\citep{goodmanbacon2021,dechaisemartin2020}. I therefore use a stacked
design \citep{cengiz2019}. Each stack $s$ consists of the firms that switch in
year $g$ to owners from origin $o$, together with clean comparison firms that
do not change ownership within the window $[g-4,g+6]$. Within a stack,
$D^s_{c}=\one[c=o_s]$ is defined for treated and comparison firms alike, so
that comparison firms supply the counterfactual path of trade with $o$. The
estimating equation is
\begin{equation}
X^s_{ickt} = \exp\Big(\sum_{e\neq-1}\beta_e\,\one[t-g_s=e]\,T^s_i\,D^s_c +
\alpha^s_{it} + \gamma^s_{ckt} + \delta^s_{ick}\Big)\,\epsilon^s_{ickt},
\label{eq:es}
\end{equation}
where $T_i^s$ indicates treated firms and all fixed effects are
stack-specific. Comparison firms are drawn from the groups in
Section~\ref{sec:data}, matched on pre-period industry, size, trading status,
and trade with $o$. As robustness, I report the estimators of
\citet{callaway2021}, \citet{sun2021}, and \citet{borusyak2024}. Event plots
report pre-event coefficients, confidence intervals, and the number of
treated firms supporting each horizon.

\subsection{Margins}
\label{sec:margins}

To map the estimates into Proposition~\ref{prop:margins}, I estimate versions
of \eqref{eq:main} and \eqref{eq:es} for each margin:
\begin{enumerate}[label=(\roman*),itemsep=0.1em]
\item \emph{Country entry}: $\one[X_{ict}>0]$ at the firm--country--year level,
in a linear probability model with firm-by-year, country-by-industry-by-year,
and firm-by-country fixed effects.
\item \emph{Product breadth}: the number of products traded with $c$, by
PPML.
\item \emph{New pairs}: trade in country--product pairs not active in the
pre-period.
\item \emph{Incumbent flows}: trade in pairs active in $t^*-1$, retaining zeros
after the transition. Conditioning on pairs that remain active afterwards
would select on an outcome of treatment and is avoided.
\end{enumerate}
The average value per active product is reported descriptively, since
changes in product composition move it even when no incumbent flow changes.

\subsection{Firm-level diversification and third countries}

The within-firm comparison absorbs any effect common to the firm's other
partners. Outcomes defined at the firm level---the number of partners, the
Herfindahl index, the non-U.S.\ share, total trade, and trade with countries
other than $o$---require a comparison across firms:
\begin{equation}
Y_{it} = \exp\big(\beta^{F}\,\Post_{it}\,T_i + \alpha_i + \kappa_{s(i)t}\big)\,u_{it},
\label{eq:firm}
\end{equation}
estimated in the stacked design, with firm fixed effects $\alpha_i$ and
industry-by-year fixed effects $\kappa_{s(i)t}$ (linear models for the
Herfindahl index and shares). Because this comparison does not difference out
firm-level shocks, it rests on a stronger assumption than \eqref{eq:main}. The
most informative comparison groups are firms sold to non-immigrant owners and
firms sold to immigrants from other origins, which share the disruption of a
sale. I interpret $\beta^{F}$ for third countries in light of
Proposition~\ref{prop:interdep}.

\subsection{Inference and power}

Standard errors are clustered by firm, the level at which treatment is
assigned \citep{abadie2023}. Because incoming owners are concentrated in a
limited number of origin countries, clustering by origin is unreliable. I
therefore complement firm-clustered inference with randomization inference
\citep{young2019}: reassigning each switching firm a placebo origin drawn from
countries with similar pre-period trade characteristics, but distant from the
true origin in language and region to avoid capturing spillovers, and
re-estimating \eqref{eq:main}. \planned{Minimum detectable effects for the
main coefficient and for each margin, computed from pre-period variance and
the number of switching firms trading in each direction, to be reported here.}

With roughly 1,000 transitions, the analysis distinguishes a small set of
pre-specified core tests---the main effect and event study by direction, the
four margins, pre-transition trading status, consumer versus intermediate
goods, and firm-level diversification---from exploratory analyses of finer
heterogeneity, spillovers to nearby countries, and reversals.

\subsection{Threats to validity}

Table~\ref{tab:threats} lists the main threats and the responses to each.

\begin{table}[htbp]\centering\footnotesize
\caption{Threats to validity}
\label{tab:threats}
\renewcommand{\arraystretch}{1.2}
\begin{tabularx}{\textwidth}{>{\raggedright\arraybackslash}p{3.2cm}YY}
\toprule
Threat & Why it matters & Response\\
\midrule
Selection on anticipated trade with $o$ & A buyer foresees growth in trade with $o$, e.g.\ a contract in prospect & Event-study leads; exclude transitions where trade with $o$ rises in $t^*-1$; compare with buyers from other origins; seller-initiated sales\\
Shocks to Canada--$o$ trade in the firm's products & A trade agreement or tariff change raises trade with $o$ at the same time & Country $\times$ product $\times$ year fixed effects\\
General effects of a sale & New owners reorganize, invest, or change suppliers & Firm $\times$ year fixed effects; transfers between non-immigrant owners as placebo\\
Network-driven matching & Firms already trading with $o$ attract buyers from $o$ & Firm $\times$ country $\times$ product fixed effects; separate results by pre-period trading status\\
Mismeasured owner identity & Early-cohort immigrants coded as non-immigrants; mixed-origin ownership & Restrict by seller age; robustness to ownership definitions and shares\\
Mechanical changes in reporting & Mergers or identifier changes alter reported trade & Require identifier continuity; exclude mergers and consolidations\\
Zeros and large traders & Log specifications drop zeros; a few firms dominate value & PPML; participation outcomes; results by firm size; unweighted linear models\\
Staggered timing & Two-way fixed effects biased under heterogeneous effects & Stacked design; robust estimators\\
\bottomrule
\end{tabularx}
\end{table}

Two placebo exercises complement the event study. First, I assign each
switching firm a placebo origin with similar pre-period trade characteristics
but no linguistic or regional proximity to the true origin; the effect should
be zero. Second, I estimate \eqref{eq:main} for firms sold to non-immigrant
owners, assigning each the origin of an immigrant buyer in the same industry
and year; the effect should again be zero.

%=====================================================================
\section{Results}
\label{sec:results}
%=====================================================================

\planned{This section is outlined. Each subsection lists the planned exhibit,
what it tests, and how alternative findings would be read.}

\subsection{The cross-sectional origin gradient}

\planned{Replicate the origin gradient of \citet{fungliu2025} for imports and
exports, among all firms and among the firms that later change hands. Split
the gradient into country participation, product breadth, and value per
relationship. This establishes the benchmark the causal estimates are compared
with: the share of the cross-sectional gradient reproduced by the
within-firm effect of a change in owner is a measure of how much of the
gradient is portable.}

\subsection{Trade with the owner's country of origin}

\figshell{Event study: trade with the incoming owner's country of origin}{%
Two panels (imports; exports). Coefficients $\beta_e$ from
equation~\eqref{eq:es}, $e\in[-4,6]$, with 95\% confidence intervals and the
number of treated firms at each horizon. Flat pre-event coefficients support
the identifying assumption; a gradual post-event rise is consistent with the
fixed-cost (search) channel, an immediate jump with the variable-cost or demand
channel.}

\begin{table}[htbp]\centering\small
\caption{Effect of a transition to immigrant ownership on trade with the owner's origin (shell)}
\label{tab:main}
\begin{tabular}{lcccc}
\toprule
 & \multicolumn{2}{c}{Imports} & \multicolumn{2}{c}{Exports}\\
\cmidrule(lr){2-3}\cmidrule(lr){4-5}
 & (1) & (2) & (3) & (4)\\
\midrule
$\Post\times\one[c=o]$ & -- & -- & -- & -- \\
 & (--) & (--) & (--) & (--) \\
\midrule
Firm $\times$ year FE & Yes & Yes & Yes & Yes \\
Country $\times$ product $\times$ year FE & Yes & Yes & Yes & Yes \\
Firm $\times$ country $\times$ product FE & Yes & Yes & Yes & Yes \\
Comparison group & All & Sold firms & All & Sold firms \\
Randomization-inference $p$-value & -- & -- & -- & -- \\
Switching firms & -- & -- & -- & -- \\
Observations & -- & -- & -- & -- \\
\bottomrule
\end{tabular}
\end{table}

\result{Main estimates by direction and comparison group.}

\subsection{Margins of adjustment}

\planned{Table with columns for country entry, product breadth, new pairs, and
incumbent flows, separately by direction. Reading: a response concentrated in
country entry and new pairs with flat incumbent flows supports the fixed-cost
channels (Table~\ref{tab:predictions}, P2--P3); proportional growth of
incumbent flows supports the variable-cost channel; both may operate.}

\subsection{Firm-level diversification}

\planned{Estimates of equation~\eqref{eq:firm} for the number of partner
countries, the number of country--product relationships, the Herfindahl index,
and the non-U.S.\ share, by direction. Test of
Proposition~\ref{prop:divers}(ii): split by whether the origin's pre-period
share is below the firm's Herfindahl index.}

\subsection{Creation or diversion}

\planned{Effects on trade with countries other than $o$ and on total trade,
across firms. Reading via Proposition~\ref{prop:interdep}: a rise in
third-country imports indicates complementarity in sourcing; a fall indicates
substitution and diversion. Report effects on trade with the United States
separately, since they matter most for national diversification.}

\subsection{Robustness}

\planned{Alternative ownership definitions (majority versus largest owner;
excluding mixed-origin owners; excluding older sellers); alternative event
windows; 2001--2019 imports; robust staggered estimators; dropping the
largest traders; linear models in levels; excluding transitions with
anticipatory trade with $o$; placebo origins; transfers between non-immigrant
owners.}

%=====================================================================
\section{Mechanisms}
\label{sec:mechanisms}
%=====================================================================

\planned{Each test interacts the treatment term with one dimension implied by
the model. The aim is to rule channels in or out, not to accumulate
heterogeneity results. Tests (1)--(4) are pre-specified; (5)--(8) are
exploratory given the sample size.}

\begin{table}[htbp]\centering\footnotesize
\caption{Mechanism tests and what each channel predicts}
\label{tab:mechtests}
\renewcommand{\arraystretch}{1.2}
\begin{tabularx}{\textwidth}{>{\raggedright\arraybackslash}p{2.4cm}>{\raggedright\arraybackslash}p{3.2cm}YYY}
\toprule
Test & Specification & Supports fixed cost if\ldots & Supports variable cost if\ldots & Supports demand if\ldots\\
\midrule
(1) Margins & Extensive vs.\ incumbent outcomes & Effect in new pairs; incumbents flat & Incumbent flows grow & Incumbent consumer imports grow\\
(2) Pre-transition status & Treatment $\times$ firm traded with $o$ before & Effect mainly among firms new to $o$ & Similar in both groups & Similar in both groups\\
(3) Direction & Imports vs.\ exports & Both directions & Both directions & Imports only\\
(4) Goods type & BEC consumer vs.\ intermediate & Both types & Both types & Consumer goods only\\
(5) Breadth & Number of products entering with $o$ & Many at once (country cost) & Few & Consumer goods only\\
(6) Information intensity & Treatment $\times$ Rauch differentiated; country friction index & Larger where frictions high & Larger where frictions high & No pattern\\
(7) Timing & Event-time profile & Gradual & Immediate & Immediate\\
(8) Spillovers & Treatment for countries near $o$ in language or region & Smaller positive effects & Smaller positive effects & None\\
\bottomrule
\end{tabularx}
\end{table}

\paragraph{Owner involvement.} \planned{If ownership shares or the owner's
employment earnings from the firm are observed, test whether effects are
concentrated among owners with controlling stakes or active roles. The
knowledge channel predicts larger effects when the owner can act on her
knowledge; passive ownership should not move trade. \verify{Availability of
ownership shares and owner earnings.}}

\paragraph{Co-national hiring.} \planned{An incoming owner may hire employees
from her origin country who carry the relevant knowledge. If employees'
immigrant status and origin can be linked, test whether the effect survives
excluding firms that hire co-national workers after the transition, or treat
hiring as a mediator. \verify{Availability of linked employee records.}}

\paragraph{Owner preferences.} \planned{A preference for home-country suppliers
based on trust rather than cost would raise trade with $o$ without lowering
measured frictions and would not be confined to consumer goods. It predicts
reallocation from other sources toward $o$ within products the firm already
imports, without a gradient in information intensity. Discuss how far the
tests above can separate it from the cost channels.}

%=====================================================================
\section{Discussion}
\label{sec:discussion}
%=====================================================================

\planned{This section is outlined.}

\paragraph{Interpreting the channel.} \planned{If the fixed-cost channel
dominates, immigrant owners diversify the firms they acquire by opening new
relationships, and the effect should persist once relationships are formed.
If the variable-cost channel dominates, owners deepen existing relationships
and the advantage is tied to their continued involvement. If the demand
channel dominates, the effect reflects the composition of the firm's
customers rather than lower trade costs, and has different welfare and
policy implications.}

\paragraph{Portability.} \planned{Compare the within-firm effect with the
cross-sectional gradient among comparable incumbent firms. The ratio is a
benchmark for how much of the immigrant--origin trade link reflects knowledge
that moves with the person, the assumption underlying quantitative
evaluations such as \citet{chikissaygin2026}. The comparison is between an
effect on acquired firms several years after a transition and a long-run
gradient, so it is a benchmark rather than an exact decomposition.}

\paragraph{From firms to the nation.} \planned{A back-of-the-envelope
aggregation: multiply the estimated effects by pre-period trade of switching
firms to obtain the contribution of ownership transitions to Canada's trade
with origin countries and to its non-U.S.\ share, and compare with the
observed change over the period. Discuss how third-country effects
(creation versus diversion, in particular from the United States) change
the national calculation.}

\paragraph{External validity.} \planned{The estimates apply to incumbent firms
acquired by immigrants, which are selected on the gains from immigrant
ownership (Proposition~\ref{prop:selection}). They should not be extrapolated
mechanically to immigrant-founded firms, to firms acquired at random, or to
changes in immigration policy.}

\paragraph{Policy.} \planned{Business succession, as owners retire, determines
who controls a large share of Canadian firms over the coming decade;
immigration programs that admit entrepreneurs, some of which allow the purchase
of existing businesses \verify{program details}, influence who the buyers are.
Discuss what the estimates imply for these programs and for trade
diversification strategies, while noting that higher trade with an origin
country is not by itself evidence of higher productivity or welfare.}

\paragraph{Limitations.} \planned{Merchandise trade only (no services); values
without prices, so price-based tests of supplier search are not available;
measurement of owner origin; limited power for fine heterogeneity and
reversals.}

%=====================================================================
\section{Conclusion}
\label{sec:conclusion}
%=====================================================================

\planned{To be written. Thoughts for the conclusion: (i) restate the question
as whether the person who owns a firm changes where it trades, and why that
matters for national diversification that is ultimately decided by firms;
(ii) summarize the headline effect, the margin, and the channel; (iii) state
the portability result relative to the cross-sectional literature; (iv) note
open questions---the long-run persistence of relationships formed under
immigrant owners, effects on firm performance (sales, employment, survival)
and on the productivity of sourcing, and the services trade that customs
records do not capture.}

%=====================================================================
\appendix
\section{Proofs}
\label{app:proofs}
%=====================================================================

\paragraph{Proof of Proposition~\ref{prop:bilateral}.} Under
Assumption~\ref{ass:sep}(i), $D_{io}$ enters only the terms of \eqref{eq:Pi}
involving $c=o$, through \eqref{eq:tau}--\eqref{eq:demand}. The optimal
choices for $c\neq o$ are unchanged by Lemma~\ref{lem:set}, so
$X_{ickd}$ is unchanged for $c\neq o$. For $c=o$, each channel weakly raises
$\Phi_{iokd}\tau_{iod}^{-\varepsilon_d}$ or weakly lowers $f_{iokd}$ or
$F_{iod}$, which weakly enlarges the active set and weakly raises the value of
every active relationship. \qed

\paragraph{Proof of Proposition~\ref{prop:margins}.} (a) With
$\lambda_f=\lambda_F=\eta=0$, \eqref{eq:tau} implies
$\Delta\ln\tau_{iod}^{-\varepsilon_d}=\varepsilon_d\lambda_\tau\iota_o$; by
Assumption~\ref{ass:sep}(ii), every relationship active before and after grows
by this amount. The threshold condition in Lemma~\ref{lem:set} is relaxed by
the same factor, so a relationship enters if and only if
$\Phi\tau^{-\varepsilon}\geq f$ after the change but not before.
(b) With only $\lambda_f>0$, $\Phi_{iokd}\tau_{iod}^{-\varepsilon_d}$ is
unchanged, so incumbent values are unchanged; a pair becomes active exactly
when its variable profit lies in the stated interval. (c) With only
$\lambda_F>0$, the product-level conditions are unchanged. If the firm was
already active with $o$, $\Omega_{iod}\geq F_{iod}$ held before and continues
to hold, and no choice changes. If it was not, it becomes active exactly when
$\Omega_{iod}$ lies in the stated interval, and it then activates every $k$
with $\Phi\tau^{-\varepsilon}\geq f$, by Lemma~\ref{lem:set}. (d) Follows from
\eqref{eq:demand}, Assumption~\ref{ass:demand}, and Assumption~\ref{ass:sep}(ii).
\qed

\paragraph{Proof of Proposition~\ref{prop:divers}.} (i) follows from
Proposition~\ref{prop:bilateral}. For (ii), let $X=\sum_c X_c$ and
$H=\sum_c X_c^2/X^2$. Raising $X_o$ with $X_c$ fixed for $c\neq o$,
\[
\frac{\partial H}{\partial X_o}=\frac{2X_oX^2-2X\sum_cX_c^2}{X^4}
=\frac{2}{X}\,\big(s_o-H\big),
\]
which is negative if and only if $s_o<H$. (iii) follows because trade with the
United States is unchanged while total trade weakly rises. \qed

\paragraph{Proof of Proposition~\ref{prop:interdep}.} In the model of
\citet{antras2017}, the firm's input purchases from $j$ equal its share
$T_j(\tau_{ij}w_j)^{-\theta}/\Theta_i$ times total input spending, which is
proportional to $\Theta_i^{(\sigma-1)/\theta}$. Hence $X_{ij}\propto
T_j(\tau_{ij}w_j)^{-\theta}\Theta_i^{(\sigma-1)/\theta-1}$. A fall in
$\ln\tau_{io}$ raises $\Theta_i$ and leaves $T_j(\tau_{ij}w_j)^{-\theta}$
unchanged for $j\neq o$, so $X_{ij}$ moves in the direction of
$(\sigma-1)/\theta-1$, in the same proportion for every $j\neq o$, proving (a).
For (b), $\ln X_{io}-\ln X_{ij}=\ln[T_o(\tau_{io}w_o)^{-\theta}]-
\ln[T_j(\tau_{ij}w_j)^{-\theta}]$, which changes by $\theta\lambda_\tau\iota_o$.
The extensive-margin statement follows the comparative statics in Section~3.2 of \citet{antras2017}: in the complements case a reduction in any $\tau_{ij}$ weakly expands the sourcing set, and (their Proposition~3) weakly raises purchases from all countries. In the knife-edge case $(\sigma-1)/\theta=1$, adding one source does not affect the decision to add any other, which is the separable baseline of Section~\ref{sec:model}. \qed

\section{Data Construction}
\label{app:data}

\planned{To be written: (i) TIC and TEC coverage, thresholds, and treatment of
re-exports; (ii) HS concordance across revisions and HS8-to-HS6 aggregation for
exports; (iii) construction of owner origin and immigrant status, and the
handling of multiple owners; (iv) identification of mergers and identifier
changes; (v) construction of the set of potential relationships; (vi) sample
restrictions and their cumulative effect on firm counts and trade coverage;
(vii) disclosure rules.}

%=====================================================================
\begin{thebibliography}{99}
\setlength{\itemsep}{0.15em}
\small

\bibitem[Abadie et~al.(2023)Abadie, Athey, Imbens, and Wooldridge]{abadie2023}
Abadie, A., S. Athey, G. W. Imbens, and J. M. Wooldridge (2023). ``When Should You Adjust Standard Errors for Clustering?'' \textit{Quarterly Journal of Economics} 138(1), 1--35.

\bibitem[Albornoz et~al.(2012)Albornoz, Calvo Pardo, Corcos, and Ornelas]{albornoz2012}
Albornoz, F., H. F. Calvo Pardo, G. Corcos, and E. Ornelas (2012). ``Sequential Exporting.'' \textit{Journal of International Economics} 88(1), 17--31.

\bibitem[Allen(2014)]{allen2014}
Allen, T. (2014). ``Information Frictions in Trade.'' \textit{Econometrica} 82(6), 2041--2083.

\bibitem[Andrews et~al.(2017)Andrews, Schank, and Upward]{andrews2017}
Andrews, M., T. Schank, and R. Upward (2017). ``Do Foreign Workers Reduce Trade Barriers? Microeconomic Evidence.'' \textit{The World Economy} 40(9), 1750--1774.

\bibitem[Antr\`as et~al.(2017)Antr\`as, Fort, and Tintelnot]{antras2017}
Antr\`as, P., T. C. Fort, and F. Tintelnot (2017). ``The Margins of Global Sourcing: Theory and Evidence from US Firms.'' \textit{American Economic Review} 107(9), 2514--2564.

\bibitem[Ariu(2022)]{ariu2022}
Ariu, A. (2022). ``Foreign Workers, Product Quality, and Trade: Evidence from a Natural Experiment.'' \textit{Journal of International Economics} 139, 103686.

\bibitem[Becker and Hvide(2022)]{becker2022}
Becker, S. O., and H. K. Hvide (2022). ``Entrepreneur Death and Startup Performance.'' \textit{Review of Finance} 26(1), 163--185.

\bibitem[Bennedsen et~al.(2007)Bennedsen, Nielsen, P\'erez-Gonz\'alez, and Wolfenzon]{bennedsen2007}
Bennedsen, M., K. M. Nielsen, F. P\'erez-Gonz\'alez, and D. Wolfenzon (2007). ``Inside the Family Firm: The Role of Families in Succession Decisions and Performance.'' \textit{Quarterly Journal of Economics} 122(2), 647--691.

\bibitem[Bernard et~al.(2007)Bernard, Jensen, Redding, and Schott]{bernard2007}
Bernard, A. B., J. B. Jensen, S. J. Redding, and P. K. Schott (2007). ``Firms in International Trade.'' \textit{Journal of Economic Perspectives} 21(3), 105--130.

\bibitem[Bernard et~al.(2011)Bernard, Redding, and Schott]{bernard2011}
Bernard, A. B., S. J. Redding, and P. K. Schott (2011). ``Multiproduct Firms and Trade Liberalization.'' \textit{Quarterly Journal of Economics} 126(3), 1271--1318.

\bibitem[Bertrand and Schoar(2003)]{bertrandschoar2003}
Bertrand, M., and A. Schoar (2003). ``Managing with Style: The Effect of Managers on Firm Policies.'' \textit{Quarterly Journal of Economics} 118(4), 1169--1208.

\bibitem[Bisztray et~al.(2018)Bisztray, Koren, and Szeidl]{bisztray2018}
Bisztray, M., M. Koren, and A. Szeidl (2018). ``Learning to Import from Your Peers.'' \textit{Journal of International Economics} 115, 242--258.

\bibitem[Bonadio(2026)]{bonadio2026}
Bonadio, B. (2026). ``Migrants, Trade, and Market Access.'' \textit{Review of Economics and Statistics} 108(2), 436--451.

\bibitem[Borusyak et~al.(2024)Borusyak, Jaravel, and Spiess]{borusyak2024}
Borusyak, K., X. Jaravel, and J. Spiess (2024). ``Revisiting Event-Study Designs: Robust and Efficient Estimation.'' \textit{Review of Economic Studies} 91(6), 3253--3285.

\bibitem[Callaway and Sant'Anna(2021)]{callaway2021}
Callaway, B., and P. H. C. Sant'Anna (2021). ``Difference-in-Differences with Multiple Time Periods.'' \textit{Journal of Econometrics} 225(2), 200--230.

\bibitem[Cardoso and Ramanarayanan(2022)]{cardoso2022}
Cardoso, M., and A. Ramanarayanan (2022). ``Immigrants and Exports: Firm-Level Evidence from Canada.'' \textit{Canadian Journal of Economics} 55(3), 1250--1293.

\bibitem[Caselli et~al.(2020)Caselli, Koren, Lisicky, and Tenreyro]{caselli2020}
Caselli, F., M. Koren, M. Lisicky, and S. Tenreyro (2020). ``Diversification Through Trade.'' \textit{Quarterly Journal of Economics} 135(1), 449--502.

\bibitem[Cengiz et~al.(2019)Cengiz, Dube, Lindner, and Zipperer]{cengiz2019}
Cengiz, D., A. Dube, A. Lindner, and B. Zipperer (2019). ``The Effect of Minimum Wages on Low-Wage Jobs.'' \textit{Quarterly Journal of Economics} 134(3), 1405--1454.

\bibitem[Chaney(2008)]{chaney2008}
Chaney, T. (2008). ``Distorted Gravity: The Intensive and Extensive Margins of International Trade.'' \textit{American Economic Review} 98(4), 1707--1721.

\bibitem[Chaney(2014)]{chaney2014}
Chaney, T. (2014). ``The Network Structure of International Trade.'' \textit{American Economic Review} 104(11), 3600--3634.

\bibitem[Chikis and Saygin(2026)]{chikissaygin2026}
Chikis, C. A., and B. Saygin (2026). ``Old Roots, New Routes: The Trade Patterns of Immigrant-Founded Firms.'' Working paper, University of Chicago.

\bibitem[Cohen et~al.(2017)Cohen, Gurun, and Malloy]{cohen2017}
Cohen, L., U. G. Gurun, and C. Malloy (2017). ``Resident Networks and Corporate Connections: Evidence from World War II Internment Camps.'' \textit{Journal of Finance} 72(1), 207--248.

\bibitem[Combes et~al.(2005)Combes, Lafourcade, and Mayer]{combes2005}
Combes, P.-P., M. Lafourcade, and T. Mayer (2005). ``The Trade-Creating Effects of Business and Social Networks: Evidence from France.'' \textit{Journal of International Economics} 66(1), 1--29.

\bibitem[Correia et~al.(2020)Correia, Guimar\~aes, and Zylkin]{correia2020}
Correia, S., P. Guimar\~aes, and T. Zylkin (2020). ``Fast Poisson Estimation with High-Dimensional Fixed Effects.'' \textit{Stata Journal} 20(1), 95--115.

\bibitem[Das et~al.(2007)Das, Roberts, and Tybout]{das2007}
Das, S., M. J. Roberts, and J. R. Tybout (2007). ``Market Entry Costs, Producer Heterogeneity, and Export Dynamics.'' \textit{Econometrica} 75(3), 837--873.

\bibitem[de~Chaisemartin and D'Haultf{\oe}uille(2020)]{dechaisemartin2020}
de Chaisemartin, C., and X. D'Haultf{\oe}uille (2020). ``Two-Way Fixed Effects Estimators with Heterogeneous Treatment Effects.'' \textit{American Economic Review} 110(9), 2964--2996.

\bibitem[Eaton et~al.(2011)Eaton, Kortum, and Kramarz]{eaton2011}
Eaton, J., S. Kortum, and F. Kramarz (2011). ``An Anatomy of International Trade: Evidence from French Firms.'' \textit{Econometrica} 79(5), 1453--1498.

\bibitem[Egger et~al.(2019)Egger, Erhardt, and Lassmann]{egger2019}
Egger, P. H., K. Erhardt, and A. Lassmann (2019). ``Immigration and Firms' Integration in International Production Networks.'' \textit{European Economic Review} 111, 1--34.

\bibitem[Erbahar and Gen\c{c}osmano\u{g}lu(2023)]{erbahar2023}
Erbahar, A., and \"O. T. Gen\c{c}osmano\u{g}lu (2023). ``Migrants and Imports: Evidence from Dutch Firms.'' \textit{Economica} 90(360), 1204--1228.

\bibitem[Fung et~al.(2019)Fung, Grekou, and Liu]{fung2019}
Fung, L., D. Grekou, and H. Liu (2019). ``The Impact of Immigrant Business Ownership on International Trade.'' Analytical Studies Branch Research Paper Series No.~426, Statistics Canada.

\bibitem[Fung and Liu(2025)]{fungliu2025}
Fung, L., and H. Liu (2025). ``Immigrant Business Ownership and Imports in Canada.'' \textit{Canadian Journal of Economics} 58(3), 961--989.

\bibitem[Goodman-Bacon(2021)]{goodmanbacon2021}
Goodman-Bacon, A. (2021). ``Difference-in-Differences with Variation in Treatment Timing.'' \textit{Journal of Econometrics} 225(2), 254--277.

\bibitem[Gould(1994)]{gould1994}
Gould, D. M. (1994). ``Immigrant Links to the Home Country: Empirical Implications for U.S. Bilateral Trade Flows.'' \textit{Review of Economics and Statistics} 76(2), 302--316.

\bibitem[Guadalupe et~al.(2012)Guadalupe, Kuzmina, and Thomas]{guadalupe2012}
Guadalupe, M., O. Kuzmina, and C. Thomas (2012). ``Innovation and Foreign Ownership.'' \textit{American Economic Review} 102(7), 3594--3627.

\bibitem[Hatzigeorgiou and Lodefalk(2016)]{hatzigeorgiou2016}
Hatzigeorgiou, A., and M. Lodefalk (2016). ``Migrants' Influence on Firm-Level Exports.'' \textit{Journal of Industry, Competition and Trade} 16(4), 477--497.

\bibitem[Head and Ries(1998)]{headries1998}
Head, K., and J. Ries (1998). ``Immigration and Trade Creation: Econometric Evidence from Canada.'' \textit{Canadian Journal of Economics} 31(1), 47--62.

\bibitem[Hiller(2013)]{hiller2013}
Hiller, S. (2013). ``Does Immigrant Employment Matter for Export Sales? Evidence from Denmark.'' \textit{Review of World Economics} 149(2), 369--394.

\bibitem[Imbs and Wacziarg(2003)]{imbs2003}
Imbs, J., and R. Wacziarg (2003). ``Stages of Diversification.'' \textit{American Economic Review} 93(1), 63--86.

\bibitem[Kaufmann et~al.(2011)Kaufmann, Kraay, and Mastruzzi]{kaufmann2011}
Kaufmann, D., A. Kraay, and M. Mastruzzi (2011). ``The Worldwide Governance Indicators: Methodology and Analytical Issues.'' \textit{Hague Journal on the Rule of Law} 3(2), 220--246.

\bibitem[Koren and Tenreyro(2007)]{koren2007}
Koren, M., and S. Tenreyro (2007). ``Volatility and Development.'' \textit{Quarterly Journal of Economics} 122(1), 243--287.

\bibitem[Lenoir et~al.(2023)Lenoir, Martin, and Mejean]{lenoir2023}
Lenoir, C., J. Martin, and I. Mejean (2023). ``Search Frictions in International Goods Markets.'' \textit{Journal of the European Economic Association} 21(1), 326--366.

\bibitem[Melitz(2003)]{melitz2003}
Melitz, M. J. (2003). ``The Impact of Trade on Intra-Industry Reallocations and Aggregate Industry Productivity.'' \textit{Econometrica} 71(6), 1695--1725.

\bibitem[Mion and Opromolla(2014)]{mionopromolla2014}
Mion, G., and L. D. Opromolla (2014). ``Managers' Mobility, Trade Performance, and Wages.'' \textit{Journal of International Economics} 94(1), 85--101.

\bibitem[Mitaritonna et~al.(2017)Mitaritonna, Orefice, and Peri]{mitaritonna2017}
Mitaritonna, C., G. Orefice, and G. Peri (2017). ``Immigrants and Firms' Outcomes: Evidence from France.'' \textit{European Economic Review} 96, 62--82.

\bibitem[Monarch and Schmidt-Eisenlohr(2023)]{monarch2023}
Monarch, R., and T. Schmidt-Eisenlohr (2023). ``Longevity and the Value of Trade Relationships.'' \textit{Journal of International Economics} 145, 103842.

\bibitem[Morales et~al.(2019)Morales, Sheu, and Zahler]{morales2019}
Morales, E., G. Sheu, and A. Zahler (2019). ``Extended Gravity.'' \textit{Review of Economic Studies} 86(6), 2668--2712.

\bibitem[Muendler and Rauch(2018)]{muendler2018}
Muendler, M.-A., and J. E. Rauch (2018). ``Do Employee Spinoffs Learn Markets from Their Parents? Evidence from International Trade.'' \textit{European Economic Review} 105, 159--173.

\bibitem[Ottaviano et~al.(2018)Ottaviano, Peri, and Wright]{ottaviano2018}
Ottaviano, G. I. P., G. Peri, and G. C. Wright (2018). ``Immigration, Trade and Productivity in Services: Evidence from U.K. Firms.'' \textit{Journal of International Economics} 112, 88--108.

\bibitem[Parsons and V\'ezina(2018)]{parsonsvezina2018}
Parsons, C., and P.-L. V\'ezina (2018). ``Migrant Networks and Trade: The Vietnamese Boat People as a Natural Experiment.'' \textit{Economic Journal} 128(612), F210--F234.

\bibitem[Patault and Lenoir(2024)]{patault2024}
Patault, B., and C. Lenoir (2024). ``Customer Capital Spillovers: Evidence from Sales Managers in International Markets.'' \textit{American Economic Journal: Applied Economics} 16(4), 404--443.

\bibitem[Peri and Requena-Silvente(2010)]{peri2010}
Peri, G., and F. Requena-Silvente (2010). ``The Trade Creation Effect of Immigrants: Evidence from the Remarkable Case of Spain.'' \textit{Canadian Journal of Economics} 43(4), 1433--1459.

\bibitem[Rauch(1999)]{rauch1999}
Rauch, J. E. (1999). ``Networks versus Markets in International Trade.'' \textit{Journal of International Economics} 48(1), 7--35.

\bibitem[Rauch(2001)]{rauch2001}
Rauch, J. E. (2001). ``Business and Social Networks in International Trade.'' \textit{Journal of Economic Literature} 39(4), 1177--1203.

\bibitem[Rauch and Trindade(2002)]{rauchtrindade2002}
Rauch, J. E., and V. Trindade (2002). ``Ethnic Chinese Networks in International Trade.'' \textit{Review of Economics and Statistics} 84(1), 116--130.

\bibitem[Santos Silva and Tenreyro(2006)]{santossilva2006}
Santos Silva, J. M. C., and S. Tenreyro (2006). ``The Log of Gravity.'' \textit{Review of Economics and Statistics} 88(4), 641--658.

\bibitem[Startz(2025)]{startz2025}
Startz, M. (2025). ``The Value of Face-to-Face: Search and Contracting Problems in Nigerian Trade.'' Working paper.

\bibitem[Steingress(2018)]{steingress2018}
Steingress, W. (2018). ``The Causal Impact of Migration on US Trade: Evidence from Political Refugees.'' \textit{Canadian Journal of Economics} 51(4), 1312--1338.

\bibitem[Steinwender(2018)]{steinwender2018}
Steinwender, C. (2018). ``Real Effects of Information Frictions: When the States and the Kingdom Became United.'' \textit{American Economic Review} 108(3), 657--696.

\bibitem[Sun and Abraham(2021)]{sun2021}
Sun, L., and S. Abraham (2021). ``Estimating Dynamic Treatment Effects in Event Studies with Heterogeneous Treatment Effects.'' \textit{Journal of Econometrics} 225(2), 175--199.

\bibitem[Wagner et~al.(2002)Wagner, Head, and Ries]{wagner2002}
Wagner, D., K. Head, and J. Ries (2002). ``Immigration and the Trade of Provinces.'' \textit{Scottish Journal of Political Economy} 49(5), 507--525.

\bibitem[Wooldridge(2023)]{wooldridge2023}
Wooldridge, J. M. (2023). ``Simple Approaches to Nonlinear Difference-in-Differences with Panel Data.'' \textit{Econometrics Journal} 26(3), C31--C66.

\bibitem[Young(2019)]{young2019}
Young, A. (2019). ``Channeling Fisher: Randomization Tests and the Statistical Insignificance of Seemingly Significant Experimental Results.'' \textit{Quarterly Journal of Economics} 134(2), 557--598.

\end{thebibliography}

\end{document}
